%%**************************************************************
%% Vorlage fuer Bachelorarbeiten (o.ä.) der DHBW
%%
%% Autor: Tobias Dreher, Yves Fischer
%% Datum: 06.07.2011
%%
%% Autor: Michael Gruben
%% Datum: 15.05.2013
%%**************************************************************

%!TEX root = ../dokumentation.tex

%
% Nahezu alle Einstellungen koennen hier getaetigt werden
%

\documentclass[%
	pdftex,
	oneside,		% Einseitiger Druck.
	12pt,			% Schriftgroesse
	parskip=half,	% Halbe Zeile Abstand zwischen Absätzen.
	headsepline,	% Linie nach Kopfzeile.
	footsepline,	% Linie vor Fusszeile.
	abstracton,	    % Abstract Überschriften
	ngerman,		% Translator
	numbers=noenddot,
]{scrreprt}

\usepackage{todonotes}

%Seitengroesse
\usepackage{fullpage}

%Zeilenumbruch und mehr
\usepackage[activate]{microtype}

% Zeichencodierung
\usepackage[utf8]{inputenc}
\usepackage[T1]{fontenc}

%For citep and citet
%\usepackage{natbib}
\usepackage[square,sort,comma,numbers]{natbib}
\usepackage{cleveref}
% Zeilenabstand
\usepackage[onehalfspacing]{setspace}

% Index-Erstellung
\usepackage{makeidx}

% Lokalisierung (neue deutsche Rechtschreibung)
\usepackage[ngerman]{babel}

% Anführungszeichen 
\usepackage[babel,german=quotes]{csquotes}
%\usepackage[style=swiss]{csquotes}


% Spezielle Tabellenform fuer Deckblatt
\usepackage{longtable}
\setlength{\tabcolsep}{10pt} %Abstand zwischen Spalten
\renewcommand{\arraystretch}{1.5} %Zeilenabstand

% Grafiken
\usepackage{graphicx}

% Mathematische Textsaetze
%\usepackage{amsmath}
%\usepackage{amssymb}

% Pakete um Textteile drehen zu können, oder eine Seite Querformat anzeigen kann.
%\usepackage{rotating}
%\usepackage{lscape}

% Farben
\usepackage{color}
\definecolor{LinkColor}{rgb}{0,0,0.2}
\definecolor{ListingBackground}{rgb}{0.92,0.92,0.92}

\newcommand{\pdftitel}{Welche Vorteile kann eine Big Data-Initiative meinem Unternehmen bringen?}
\newcommand{\autor}{Wenja Shan, Kristin Meiswinkel}
\newcommand{\arbeit}{Integrationsseminar}

% Titel, Autor und Datum
\title{\titel}
\author{\autor}
\date{\datum}

% PDF Einstellungen
\usepackage[%
	pdftitle={\pdftitel},
	pdfauthor={\autor},
	pdfsubject={\arbeit},
	pdfcreator={pdflatex, LaTeX with KOMA-Script},
	pdfpagemode=UseOutlines, % Beim Oeffnen Inhaltsverzeichnis anzeigen
	pdfdisplaydoctitle=true, % Dokumenttitel statt Dateiname anzeigen.
	pdflang=de % Sprache des Dokuments.
]{hyperref} 

% (Farb-)einstellungen für die Links im PDF
\hypersetup{%
	colorlinks=true, % Aktivieren von farbigen Links im Dokument
	linkcolor=LinkColor, % Farbe festlegen
	citecolor=LinkColor,
	filecolor=LinkColor,
	menucolor=LinkColor,
	urlcolor=LinkColor,
	bookmarksnumbered=true % Überschriftsnummerierung im PDF Inhalt anzeigen.
}

% Literaturverweise nach Harvard (mit deutschem und)
%\usepackage[dcucite]{harvard}
%\renewcommand{\harvardand}{und}

% Verschiedene Schriftarten
%\usepackage{goudysans}
%\usepackage{lmodern}
%\usepackage{libertine}
\usepackage{palatino} 

% Hurenkinder und Schusterjungen verhindern
% http://projekte.dante.de/DanteFAQ/Silbentrennung
\clubpenalty=10000
\widowpenalty=10000
\displaywidowpenalty=10000

% Quellcode
\usepackage{listings}
\lstloadlanguages{Java}
\lstset{%
	language=PHP,		 	 % Sprache des Quellcodes
	%numbers=left,           % Zelennummern links
	stepnumber=1,            % Jede Zeile nummerieren.
	numbersep=5pt,           % 5pt Abstand zum Quellcode
	numberstyle=\tiny,       % Zeichengrösse 'tiny' für die Nummern.
	breaklines=true,         % Zeilen umbrechen wenn notwendig.
	breakautoindent=true,    % Nach dem Zeilenumbruch Zeile einrücken.
	postbreak=\space,        % Bei Leerzeichen umbrechen.
	tabsize=2,               % Tabulatorgrösse 2
	basicstyle=\ttfamily\scriptsize, % Nichtproportionale Schrift, klein für den Quellcode
	showspaces=false,        % Leerzeichen nicht anzeigen.
	showstringspaces=false,  % Leerzeichen auch in Strings ('') nicht anzeigen.
	extendedchars=true,      % Alle Zeichen vom Latin1 Zeichensatz anzeigen.
	captionpos=b,            % sets the caption-position to bottom
	backgroundcolor=\color{ListingBackground} % Hintergrundfarbe des Quellcodes setzen.
}

\usepackage{caption}

% Glossar
\usepackage[
	nonumberlist, %keine Seitenzahlen anzeigen
	%acronym,      %ein Abkürzungsverzeichnis erstellen
	%section,     %im Inhaltsverzeichnis auf section-Ebene erscheinen
	toc,          %Einträge im Inhaltsverzeichnis
]{glossaries}

%Akronyme
\usepackage[printonlyused]{acronym}

% Fussnoten
\usepackage[perpage, hang, multiple, stable]{footmisc}


%Bildpfad
\graphicspath{{images/}}

%nur ein latex-Durchlauf für die Aktualisierung von Verzeichnissen nötig
\usepackage{bookmark}

%Gleitumgebungen (Bilder, Tabellen, usw\ldots) lassen sich mit H an genau der
% definierten Stelle platzieren
\usepackage{float}

% für die vertikale Platzierung von Text in Tabellen
\usepackage{array}

% für die Darstellung des Euro-Symbols
\usepackage[right]{eurosym}

% für textumflossene Grafiken
\usepackage{wrapfig}

\usepackage{longtable}
\usepackage{colortbl}

% eine Kommentarumgebung "k" (Handhabe mit \begin{k}<Kommentartext>\end{k},
% Kommentare werden rot gedruckt). Wird \% vor excludecomment{k} entfernt,
% werden keine Kommentare mehr gedruckt.
\usepackage{comment}
\specialcomment{k}{\begingroup\color{red}}{\endgroup}
%\excludecomment{k}

\usepackage{pdfpages}

\usepackage{subfigure}

 \usepackage{bibgerm}

% Ab jetzt können auch Umlaute verwendet werden

%falls pdftitel = titel der Arbeit
\newcommand{\titel}{\pdftitel}
%bei unterschiedlichen Titeln
%\newcommand{\titel}{I haben wir einen zweizeiligen
% Bachelorthesistitel}
\newcommand{\martrikelnr}{2929412}
\newcommand{\kurs}{WWI12B5}
\newcommand{\datumAbgabe}{Juli 2015}
\newcommand{\firma}{SAP SE}
\newcommand{\firmenort}{Walldorf}
\newcommand{\abgabeort}{Karlsruhe}
\newcommand{\abschluss}{Bachelor of Science}
\newcommand{\studiengang}{Wirtschaftsinformatik}
\newcommand{\dhbw}{Karlsruhe}
\newcommand{\betreuer}{Dr. Bernd Limberger}
\newcommand{\gutachter}{Torsten Rösen}
\newcommand{\zeitraum}{13.05.2015 bis 01.07.2015}
\newcommand{\arbeitsart}{\arbeit}

\makeglossaries
%!TEX root = ../dokumentation.tex

%
% vorher in Konsole folgendes aufrufen: 
%	makeglossaries makeglossaries dokumentation.acn && makeglossaries dokumentation.glo
%

%
% Glossareintraege --> referenz, name, beschreibung
% Aufruf mit \gls{...}
%
\newglossaryentry{Glossareintrag}{name={Glossareintrag},plural={Glossareinträge},description={Ein Glossar beschreibt verschiedenste Dinge in kurzen Worten}}


\begin{document}



	% Deckblatt
	\begin{spacing}{1}
		%!TEX root = ../dokumentation.tex

\begin{titlepage}
	\begin{center} 
	  {\includegraphics[height=2.6cm]{images/dhbw.png}}
	\end{center}
	\enlargethispage{20mm}
	\begin{center}
	  \vspace*{12mm}	{\LARGE\bf \titel }\\
	  \vspace*{12mm}	{\large\bf \arbeit}\\
	  \vspace*{12mm}	für die Prüfung zum\\
	  \vspace*{3mm} 	{\bf \abschluss}\\
	  \vspace*{12mm}	des Studiengangs \studiengang\\
	  \vspace*{3mm} 	an der Dualen Hochschule Baden-Württemberg \dhbw\\
	  \vspace*{12mm}	von\\
	  \vspace*{3mm} 	{\large\bf \autor}\\
	  \vspace*{12mm}	\datumAbgabe\\
	\end{center}
	\vfill
	\begin{spacing}{1.2}
	\begin{tabbing}
		mmmmmmmmmmmmmmmmmmmmmmmmmm     \= \kill
		\textbf{Bearbeitungszeitraum}  \>  \zeitraum\\
		\textbf{Matrikelnummer, Kurs}  \>  \martrikelnr, \kurs\\
		\textbf{Ausbildungsfirma}      \>  \firma, \firmenort\\
		\textbf{Betreuer}              \>  \betreuer\\
		\textbf{Gutachter}             \>  \gutachter
	\end{tabbing}
	\end{spacing}
\end{titlepage}

	\end{spacing}
	\newpage

	\renewcommand{\thepage}{\Roman{page}}
	\setcounter{page}{1}

	% Sperrvermerk
	%!TEX root = ../dokumentation.tex

\thispagestyle{empty}
% Sperrvermerk direkt hinter Titelseite
\section*{Sperrvermerk}

\vspace*{2em}

Die vorliegende Arbeit mit dem Titel {\itshape \titel} ist mit einem Sperrvermerk versehen und wird ausschließlich zu Prüfungszwecken am Studiengang {\studiengang} der Dualen Hochschule Baden-Württemberg {\abgabeort} vorgelegt.
Jede Einsichtnahme und Veröffentlichung – auch von Teilen der Arbeit – bedarf der vorherigen Zustimmung durch die {\firma}.

	\newpage
	
	% Erklärung
	%!TEX root = ../dokumentation.tex

\thispagestyle{empty}

\section*{Eidesstattliche Erklärung}
% http://www.se.dhbw-mannheim.de/fileadmin/ms/wi/dl_swm/dhbw-ma-wi-organisation-bewertung-bachelorarbeit-v2-00.pdf
\vspace*{2em}

Ich erkläre hiermit eidesstattlich, dass ich die vorliegende Arbeit selbstständig und ohne Benutzung anderer als der angegebenen Hilfsmittel angefertigt habe. Aus den benutzten Quellen direkt oder indirekt übernommene Gedanken habe ich als solche kenntlich gemacht.\\
\newline
Diese Arbeit wurde bisher in gleicher oder ähnlicher Form oder auszugsweise noch keiner anderen Prüfungsbehörde vorgelegt und auch nicht veröffentlicht.

\vspace{3em}

\abgabeort, \datumAbgabe
\vspace{4em}

\autor

	\newpage


	% Abstract
	%%!TEX root = ../dokumentation.tex

\pagestyle{empty}

\renewcommand{\abstractname}{Zusammenfassung}
\begin{abstract}
Ein Abstract ist eine prägnante Inhaltsangabe, ein Abriss ohne
Interpretation und Wertung einer wissenschaftlichen Arbeit. In DIN
1426 wird das (oder auch der) Abstract als Kurzreferat zur
Inhaltsangabe beschrieben.

\begin{description}
\item[Objektivität] soll sich jeder persönlichen Wertung enthalten
\item[Kürze] soll so kurz wie möglich sein
\item[Genauigkeit] soll genau die Inhalte und die Meinung der Originalarbeit wiedergeben
\end{description}

Üblicherweise müssen wissenschaftliche Artikel einen Abstract
enthalten, typischerweise von 100-150 Wörtern, ohne Bilder und
Literaturzitate und in einem Absatz.

Quelle \url{http://de.wikipedia.org/wiki/Abstract} Abgerufen 07.07.2011
\end{abstract}


\renewcommand{\abstractname}{Summary}
\begin{abstract}
An abstract is a brief summary of a research article, thesis, review,
conference proceeding or any in-depth analysis of a particular subject
or discipline, and is often used to help the reader quickly ascertain
the paper's purpose. When used, an abstract always appears at the
beginning of a manuscript, acting as the point-of-entry for any given
scientific paper or patent application. Abstracting and indexing
services for various academic disciplines are aimed at compiling a
body of literature for that particular subject.

The terms précis or synopsis are used in some publications to refer to
the same thing that other publications might call an "abstract". In
management reports, an executive summary usually contains more
information (and often more sensitive information) than the abstract
does.

Quelle: \url{http://en.wikipedia.org/wiki/Abstract_(summary)}

\end{abstract}

	%\newpage

	\pagestyle{plain}

	% Inhaltsverzeichnis
	\begin{spacing}{1.1}
		\setcounter{tocdepth}{2}
		%für die Anzeige von Unterkapiteln im Inhaltsverzeichnis
		%\setcounter{tocdepth}{2}
		\tableofcontents
	\end{spacing}
	\newpage
	
		% Abkürzungsverzeichnis
		\cleardoublepage
		\phantomsection \label{listofacs}
		\addcontentsline{toc}{chapter}{Abkürzungsverzeichnis}
		%!TEX root = ../dokumentation.tex

\chapter*{Abkürzungsverzeichnis}
%nur verwendete Akronyme werden letztlich im Dokument angezeigt
\begin{acronym}[YTMMM]
\setlength{\itemsep}{-\parsep}

\end{acronym}

		
		
		% Abbildungsverzeichnis
		\cleardoublepage
		\phantomsection \label{listoffig}
		\addcontentsline{toc}{chapter}{Abbildungsverzeichnis}
		\listoffigures
	
		%Tabellenverzeichnis
		\cleardoublepage
		\phantomsection \label{listoftab}
		\addcontentsline{toc}{chapter}{Tabellenverzeichnis}
		\listoftables

    \cleardoublepage
	\renewcommand{\thepage}{\arabic{page}}
	\setcounter{page}{1}
	
	% Inhalt
     %!TEX root = ../dokumentation.tex
%\todo{Beteuer fragen wegen Zitierweise (Referenzen nach oder vor dem Punkt, wenn sich diese Referenz nicht nur auf einen Satz bezieht), Fußnoten von Acronyms, Schaubilder auf englisch, Schaubilder horizontal }

\chapter{Einführung}


     \chapter{Grundlagen}


     \chapter{Planung und Analyse}

     \chapter{Konzeption des Beratungsleitfadens}


     
\chapter{Kritische Analyse des Leitfadens}

     \chapter{Fazit und Ausblick}

     	
	
	% Anhang
	\clearpage
	\renewcommand{\thepage}{\Roman{page}}
		\setcounter{page}{8}
	%\pagenumbering{roman}
	\appendix
	\chapter{Anhang}

	



	% Quellcodeverzeichnis
	%\cleardoublepage
	%\phantomsection \label{listoflist}
	%\addcontentsline{toc}{chapter}{Listings}
	%\lstlistoflistings


	% Literaturverzeichnis
	\cleardoublepage
	\phantomsection \label{listoflit}
	\addcontentsline{toc}{chapter}{Literaturverzeichnis}
	
	\nocite{*}
	%Bib style
	%\bibliographystyle{agsm} %Havard  
	%\bibliographystyle{amsplain} %Durchnummeriert
	%\bibliographystyle{amsalpha} %Kürzel für Autor und Jahr,
	\bibliographystyle{geralpha} % was Jan verwendet: alpha
	%see more: http://amath.colorado.edu/documentation/LaTeX/reference/faq/bibstyles.pdf
	
  	\bibliography{ArbeitBib}
     
   %cd 
   \cleardoublepage
   \phantomsection
   \addcontentsline{toc}{chapter}{Beigabenverzeichnis}
   \chapter*{Beigabenverzeichnis}

\begin{figure}[H] 
  \centering
     \includegraphics[scale=0.35]{cd.pdf}
\end{figure}


Im Folgenden ist die Verzeichnisstruktur der beigelegten CD dargestellt.
 
\begin{description}
  \item[Bachelorarbeit]
    Beinhaltet die Bachelorarbeit als \textit{PDF}.
    Der zugehörige \emph{\LaTeX{}} Source Code befindet sich im Unterverzeichnis \textit{LaTeX}.
    Im Unterverzeichnis \textit{DOC} befindet sich die von \textit{PDF} in \textit{DOC} konvertierte Version dieser Arbeit.
  \item[Elektronische Quellen]
  %\todo{\BibTeX funktioniert nicht}
    Beinhaltet das Literaturverzeichnis im \textit{BibTeX} Format, sowie die digital verfügbaren Quellen als \textit{PDF} oder \textit{MHT}.
\end{description}

     
     
	% Glossar
	\printglossary[style=altlist,title=Glossar]
\end{document}
